% generated by GAPDoc2LaTeX from XML source (Frank Luebeck)
\documentclass[a4paper,11pt]{report}

\usepackage[top=37mm,bottom=37mm,left=27mm,right=27mm]{geometry}
\sloppy
\pagestyle{myheadings}
\usepackage{amssymb}
\usepackage[utf8]{inputenc}
\usepackage{makeidx}
\makeindex
\usepackage{color}
\definecolor{FireBrick}{rgb}{0.5812,0.0074,0.0083}
\definecolor{RoyalBlue}{rgb}{0.0236,0.0894,0.6179}
\definecolor{RoyalGreen}{rgb}{0.0236,0.6179,0.0894}
\definecolor{RoyalRed}{rgb}{0.6179,0.0236,0.0894}
\definecolor{LightBlue}{rgb}{0.8544,0.9511,1.0000}
\definecolor{Black}{rgb}{0.0,0.0,0.0}

\definecolor{linkColor}{rgb}{0.0,0.0,0.554}
\definecolor{citeColor}{rgb}{0.0,0.0,0.554}
\definecolor{fileColor}{rgb}{0.0,0.0,0.554}
\definecolor{urlColor}{rgb}{0.0,0.0,0.554}
\definecolor{promptColor}{rgb}{0.0,0.0,0.589}
\definecolor{brkpromptColor}{rgb}{0.589,0.0,0.0}
\definecolor{gapinputColor}{rgb}{0.589,0.0,0.0}
\definecolor{gapoutputColor}{rgb}{0.0,0.0,0.0}

%%  for a long time these were red and blue by default,
%%  now black, but keep variables to overwrite
\definecolor{FuncColor}{rgb}{0.0,0.0,0.0}
%% strange name because of pdflatex bug:
\definecolor{Chapter }{rgb}{0.0,0.0,0.0}
\definecolor{DarkOlive}{rgb}{0.1047,0.2412,0.0064}


\usepackage{fancyvrb}

\usepackage{mathptmx,helvet}
\usepackage[T1]{fontenc}
\usepackage{textcomp}


\usepackage[
            pdftex=true,
            bookmarks=true,        
            a4paper=true,
            pdftitle={Written with GAPDoc},
            pdfcreator={LaTeX with hyperref package / GAPDoc},
            colorlinks=true,
            backref=page,
            breaklinks=true,
            linkcolor=linkColor,
            citecolor=citeColor,
            filecolor=fileColor,
            urlcolor=urlColor,
            pdfpagemode={UseNone}, 
           ]{hyperref}

\newcommand{\maintitlesize}{\fontsize{50}{55}\selectfont}

% write page numbers to a .pnr log file for online help
\newwrite\pagenrlog
\immediate\openout\pagenrlog =\jobname.pnr
\immediate\write\pagenrlog{PAGENRS := [}
\newcommand{\logpage}[1]{\protect\write\pagenrlog{#1, \thepage,}}
%% were never documented, give conflicts with some additional packages

\newcommand{\GAP}{\textsf{GAP}}

%% nicer description environments, allows long labels
\usepackage{enumitem}
\setdescription{style=nextline}

%% depth of toc
\setcounter{tocdepth}{1}





%% command for ColorPrompt style examples
\newcommand{\gapprompt}[1]{\color{promptColor}{\bfseries #1}}
\newcommand{\gapbrkprompt}[1]{\color{brkpromptColor}{\bfseries #1}}
\newcommand{\gapinput}[1]{\color{gapinputColor}{#1}}


\begin{document}

\logpage{[ 0, 0, 0 ]}
\begin{titlepage}
\mbox{}\vfill

\begin{center}{\maintitlesize \textbf{ typeset \mbox{}}}\\
\vfill

\hypersetup{pdftitle= typeset }
\markright{\scriptsize \mbox{}\hfill  typeset  \hfill\mbox{}}
{\Huge \textbf{ Automatic typesetting framework for common \textsf{GAP} objects, with LaTeX generation \mbox{}}}\\
\vfill

{\Huge  1.2.4 \mbox{}}\\[1cm]
{ 28 April 2026 \mbox{}}\\[1cm]
\mbox{}\\[2cm]
{\Large \textbf{ Zachariah Newbery\\
   \mbox{}}}\\
\hypersetup{pdfauthor= Zachariah Newbery\\
   }
\end{center}\vfill

\mbox{}\\
{\mbox{}\\
\small \noindent \textbf{ Zachariah Newbery\\
   }  Email: \href{mailto://me@zachnewbery.com} {\texttt{me@zachnewbery.com}}\\
  Homepage: \href{https://zachnewbery.com} {\texttt{https://zachnewbery.com}}}\\
\end{titlepage}

\newpage\setcounter{page}{2}
\newpage

\def\contentsname{Contents\logpage{[ 0, 0, 1 ]}}

\tableofcontents
\newpage

     
\chapter{\textcolor{Chapter }{Introduction}}\label{Chapter_Introduction}
\logpage{[ 1, 0, 0 ]}
\hyperdef{L}{X7DFB63A97E67C0A1}{}
{
  

 
\section{\textcolor{Chapter }{Core Framework Functions}}\label{Chapter_Introduction_Section_Core_Framework_Functions}
\logpage{[ 1, 1, 0 ]}
\hyperdef{L}{X7F60CF8C7C27BFFF}{}
{
  

 \textsf{typeset} is a package that implements a typesetting framework that can be implemented
for numerous typesetting languages as a standardised way to generate
renderable strings. 

 At it's core, it implements the function \texttt{Typeset} (\ref{Typeset}) which makes use of typesetting language\texttt{\symbol{45}}specific functions
to generate format strings. These strings are then populated with a list of
the semantic features of the \textsf{GAP} objects they represent, which is obtained from the operation \texttt{GenArgs} (\ref{GenArgs:for IsObject}). 

 An example implementation of this framework is also provided by this package
for LaTeX typesetting within chapter \ref{Chapter_LaTeX_Generation}. 

 Guidelines for extending the framework to support more types, or for
implementing the framework for another typesetting language can be found
within the contributing guidelines in the GitHub repository. 

 

\subsection{\textcolor{Chapter }{InfoTypeset}}
\logpage{[ 1, 1, 1 ]}\nobreak
\hyperdef{L}{X7A7D88777B9A0588}{}
{\noindent\textcolor{FuncColor}{$\triangleright$\enspace\texttt{InfoTypeset\index{InfoTypeset@\texttt{InfoTypeset}}
\label{InfoTypeset}
}\hfill{\scriptsize (info class)}}\\


 Info class for the \textsf{typeset} package. Set this to the following levels for different levels of information: 
\begin{itemize}
\item  0 \texttt{\symbol{45}} No messages 
\item  1 \texttt{\symbol{45}} Problems only: messages describing what went wrong,
with no messages if an operation is successful 
\item  2 \texttt{\symbol{45}} Required preamble packages: displays informations about
any required LaTeX packages that need to be added to the preamble to be
rendered. 
\item  3 \texttt{\symbol{45}} Progress: also shows
step\texttt{\symbol{45}}by\texttt{\symbol{45}}step progress of operations 
\end{itemize}
 

 Set this using, for example \texttt{SetInfoLevel(InfoTypeset, 1)}. Default value is 2. }

 

\subsection{\textcolor{Chapter }{Typeset}}
\logpage{[ 1, 1, 2 ]}\nobreak
\hyperdef{L}{X8739447F7A4111CC}{}
{\noindent\textcolor{FuncColor}{$\triangleright$\enspace\texttt{Typeset({\mdseries\slshape obj[, options]})\index{Typeset@\texttt{Typeset}}
\label{Typeset}
}\hfill{\scriptsize (function)}}\\
\textbf{\indent Returns:\ }
 A String, if \texttt{ReturnStr} option is set to \texttt{true} 



 Generates a mark\texttt{\symbol{45}}up string representing the object \mbox{\texttt{\mdseries\slshape obj}} in the given mark\texttt{\symbol{45}}up language. \textsf{GAP} options can also be added to modify the result: 
\begin{itemize}
\item  \texttt{ReturnStr} : Whether the method should return a string (\texttt{true}), or simply print the result (\texttt{false}). (default \texttt{\symbol{45}} \texttt{false}) 
\item  \texttt{LDelim} : Left Delimiter for matrices. (default \texttt{\symbol{45}} \texttt{"("}) 
\item  \texttt{RDelim} : Right Delimiter for matrices. (default \texttt{\symbol{45}} \texttt{")"}) 
\item  \texttt{Lang} : Markup language of output, currently only \texttt{"latex"} is supported. (default \texttt{\symbol{45}} \texttt{"latex"}) 
\item  \texttt{DigraphOut} : Typesetting method for Digraphs, one of \texttt{"dot"} to use raw dot within TeX, or \texttt{"dot2tex"} to convert the dot to native TeX. (default \texttt{\symbol{45}} \texttt{"dot"}) 
\item  \texttt{SubCallOpts} : Alternate \textsf{GAP} options for nested sub\texttt{\symbol{45}}objects, via a record with the same
options as the parent (but different values), or \texttt{false} if all options are to stay the same between sub\texttt{\symbol{45}}calls.
Options merging is handled by \texttt{MergeSubOptions} (\ref{MergeSubOptions}). (default \texttt{\symbol{45}} \texttt{false}) either by specifying each options as an individual GAP options like below: 
\end{itemize}
 

 
\begin{Verbatim}[commandchars=!@|,fontsize=\small,frame=single,label=Example]
  !gapprompt@gap>| !gapinput@Typeset([[1, 2], [2, 1]] : LDelim := "[", ReturnStr := true);|
  "\\left[\\begin{array}{rr}\n1 & 2 \\\\\n2 & 1 \\\\\n\\end{array}\\right)\n"
\end{Verbatim}
 

 or wrapping them in a record under an \texttt{options} GAP option, like: 

 
\begin{Verbatim}[commandchars=!@|,fontsize=\small,frame=single,label=Example]
  !gapprompt@gap>| !gapinput@Typeset([[1, 2], [2, 1]] : options := rec(LDelim := "[", ReturnStr := true));|
  "\\left[\\begin{array}{rr}\n1 & 2 \\\\\n2 & 1 \\\\\n\\end{array}\\right)\n"
\end{Verbatim}
 

 or even simply passing a record object as the optional second argument: 

 
\begin{Verbatim}[commandchars=!@|,fontsize=\small,frame=single,label=Example]
  !gapprompt@gap>| !gapinput@Typeset([[1, 2], [2, 1]], rec(LDelim := "[", ReturnStr := true));|
  "\\left[\\begin{array}{rr}\n1 & 2 \\\\\n2 & 1 \\\\\n\\end{array}\\right)\n"
\end{Verbatim}
 

 }

 

\subsection{\textcolor{Chapter }{TypesetInternal}}
\logpage{[ 1, 1, 3 ]}\nobreak
\hyperdef{L}{X860735DF791563F0}{}
{\noindent\textcolor{FuncColor}{$\triangleright$\enspace\texttt{TypesetInternal({\mdseries\slshape obj})\index{TypesetInternal@\texttt{TypesetInternal}}
\label{TypesetInternal}
}\hfill{\scriptsize (function)}}\\
\textbf{\indent Returns:\ }
 A String 



 Generates a string representation of a passed \textsf{GAP} object \mbox{\texttt{\mdseries\slshape obj}} that can be rendered by a typesetter. Called from the
top\texttt{\symbol{45}}level method \texttt{Typeset} (\ref{Typeset}), which also passes a constructed options record as the \texttt{options} \textsf{GAP} option. 

 }

 }

 
\section{\textcolor{Chapter }{Core Operations}}\label{Chapter_Introduction_Section_Core_Operations}
\logpage{[ 1, 2, 0 ]}
\hyperdef{L}{X79FAE10A87238E4C}{}
{
  

\subsection{\textcolor{Chapter }{GenArgs (for IsObject)}}
\logpage{[ 1, 2, 1 ]}\nobreak
\hyperdef{L}{X860806A279FED206}{}
{\noindent\textcolor{FuncColor}{$\triangleright$\enspace\texttt{GenArgs({\mdseries\slshape obj})\index{GenArgs@\texttt{GenArgs}!for IsObject}
\label{GenArgs:for IsObject}
}\hfill{\scriptsize (operation)}}\\
\textbf{\indent Returns:\ }
 A List of Strings 



 Generates the arguments describing the semantic definition of the passed \textsf{GAP} object \mbox{\texttt{\mdseries\slshape obj}}. This returns a list that can be used to populate a format string in any
mark\texttt{\symbol{45}}up language. If no method is installed for a type, it
will fallback to returning the list [ ViewString(obj) ]. 

 }

 }

 
\section{\textcolor{Chapter }{Constants and Utility Functions}}\label{Chapter_Introduction_Section_Constants_and_Utility_Functions}
\logpage{[ 1, 3, 0 ]}
\hyperdef{L}{X7F753FFF815225A2}{}
{
  

 

\subsection{\textcolor{Chapter }{MergeSubOptions}}
\logpage{[ 1, 3, 1 ]}\nobreak
\hyperdef{L}{X8588DF877EC7F955}{}
{\noindent\textcolor{FuncColor}{$\triangleright$\enspace\texttt{MergeSubOptions({\mdseries\slshape opts})\index{MergeSubOptions@\texttt{MergeSubOptions}}
\label{MergeSubOptions}
}\hfill{\scriptsize (function)}}\\
\textbf{\indent Returns:\ }
 A Record 



 Merges the passed options record \mbox{\texttt{\mdseries\slshape opts}} to change any values that are set in the \textsf{GAP} option \texttt{SubCallOpts}. If this option is not false (default), it can contain a record of any \textsf{GAP} options that can be passed to \texttt{Typeset} (\ref{Typeset}) which should differ for sub\texttt{\symbol{45}}calls. 

 For example, to alter the delimiters for nested objects so that the outer
object is delimited by square braces and the inner object by parentheses, the
following can be set: 

 
\begin{Verbatim}[commandchars=!@|,fontsize=\small,frame=single,label=Example]
  !gapprompt@gap>| !gapinput@MergeSubOptions(rec(ReturnStr := false, Lang := "latex", DigraphOut := "dot", RDelim := "]", LDelim := "[", SubCallOpts := rec(RDelim := ")", LDelim := "(")));|
  rec(ReturnStr := false, Lang := "latex", DigraphOut := "dot", RDelim := ")", LDelim := "(", SubCallOpts := false)
\end{Verbatim}
 

 It should be noted that \texttt{SubCallOpts} only changes the options for one level of recursion (i.e. it is set back to
the default of \texttt{false} once this function is called). To change options for more recursion levels,
the \texttt{SubCallOpts} option can be nested as many times as necessary. 

 }

 

\subsection{\textcolor{Chapter }{DEFAULT{\textunderscore}TYPESET{\textunderscore}OPTIONS}}
\logpage{[ 1, 3, 2 ]}\nobreak
\hyperdef{L}{X83056BB487424F07}{}
{\noindent\textcolor{FuncColor}{$\triangleright$\enspace\texttt{DEFAULT{\textunderscore}TYPESET{\textunderscore}OPTIONS\index{DEFAULT{\textunderscore}TYPESET{\textunderscore}OPTIONS@\texttt{DEF}\-\texttt{A}\-\texttt{U}\-\texttt{L}\-\texttt{T{\textunderscore}}\-\texttt{T}\-\texttt{Y}\-\texttt{P}\-\texttt{E}\-\texttt{S}\-\texttt{E}\-\texttt{T{\textunderscore}}\-\texttt{O}\-\texttt{P}\-\texttt{T}\-\texttt{IONS}}
\label{DEFAULTuScoreTYPESETuScoreOPTIONS}
}\hfill{\scriptsize (global variable)}}\\


 Default options record passed to \texttt{Typeset} (\ref{Typeset}). Merged with user\texttt{\symbol{45}}provided options to ensure correct
construction of options for sub\texttt{\symbol{45}}calls, whilst also allowing
option\texttt{\symbol{45}}less calls to the method. }

 }

 }

   
\chapter{\textcolor{Chapter }{LaTeX Generation}}\label{Chapter_LaTeX_Generation}
\logpage{[ 2, 0, 0 ]}
\hyperdef{L}{X877F82D67F8A4C37}{}
{
  

 
\section{\textcolor{Chapter }{LaTeX Generation Functions for GAP Objects}}\label{Chapter_LaTeX_Generation_Section_LaTeX_Generation_Functions_for_GAP_Objects}
\logpage{[ 2, 1, 0 ]}
\hyperdef{L}{X87F6833C7F252C79}{}
{
  

 The introduction in \ref{Chapter_Introduction_Section_Core_Framework_Functions} describes a powerful framework created in this package allowing for a
standardised methodology to generate typesetting strings via the semantic
features of objects. 

 This section describes the implementation of the framework for LaTeX,
providing both the invaluable functionality to typeset a subset of \textsf{GAP} objects as LaTeX strings whilst also serving as an example of how the
framework can be used for other typesetting languages. 

 It also provides some insight into the kinds of functions that would be
expected from an implementation for a different language. A
bare\texttt{\symbol{45}}boned example implementation for MathML is provided in \ref{Chapter_MathML_Example_Generation}. 

 Currently, the following types have explicit methods installed for LaTeX
generation: 

 
\begin{itemize}
\item  Rationals (Integers and Fractions) 
\item  Infinity and Negative Infinity 
\item  Internal Finite\texttt{\symbol{45}}Field Elements 
\item  Permutations 
\item  Lists 
\item  Matrices 
\item  Polynomials and Non\texttt{\symbol{45}}Polynomial Rational Functions 
\item  Character Tables 
\item  Generators for FP, PC, Matrix and Permutation Groups 
\item  Associative Words in Letter Representation 
\end{itemize}
 

 It should also be noted that \texttt{Typeset} (\ref{Typeset}) does fallback to the core library function \texttt{ViewString} (\textbf{Reference: ViewString}), which can be used for any types which do not require
LaTeX\texttt{\symbol{45}}specific formatting to be renderable in math mode
(e.g. floats). 

\subsection{\textcolor{Chapter }{GenLatexTmpl (for IsObject)}}
\logpage{[ 2, 1, 1 ]}\nobreak
\hyperdef{L}{X817F73037F401EBE}{}
{\noindent\textcolor{FuncColor}{$\triangleright$\enspace\texttt{GenLatexTmpl({\mdseries\slshape obj})\index{GenLatexTmpl@\texttt{GenLatexTmpl}!for IsObject}
\label{GenLatexTmpl:for IsObject}
}\hfill{\scriptsize (operation)}}\\
\textbf{\indent Returns:\ }
 An Unpopulated LaTeX Format String 



 Generates a format string that represents the structural definition of the
given \textsf{GAP} object \mbox{\texttt{\mdseries\slshape obj}} in LaTeX. It contains no parameter values, and will need to be populated with
the arguments representing the semantic values of the object, generated via \texttt{GenArgs} (\ref{GenArgs:for IsObject}), before it can be rendered in a LaTeX environment. 

 }

 

\subsection{\textcolor{Chapter }{CtblEntryLatex}}
\logpage{[ 2, 1, 2 ]}\nobreak
\hyperdef{L}{X793E144084AFEBCB}{}
{\noindent\textcolor{FuncColor}{$\triangleright$\enspace\texttt{CtblEntryLatex({\mdseries\slshape s})\index{CtblEntryLatex@\texttt{CtblEntryLatex}}
\label{CtblEntryLatex}
}\hfill{\scriptsize (function)}}\\
\textbf{\indent Returns:\ }
 A String 



 Formats a string representation of an entry \mbox{\texttt{\mdseries\slshape s}} returned by the undocumented function \texttt{CharacterTableDisplayStringEntryDefault} to include the LaTeX\texttt{\symbol{45}}specific bar environment for complex
conjugates. 

 }

 

\subsection{\textcolor{Chapter }{CtblLegendLatex}}
\logpage{[ 2, 1, 3 ]}\nobreak
\hyperdef{L}{X79D99604858275AF}{}
{\noindent\textcolor{FuncColor}{$\triangleright$\enspace\texttt{CtblLegendLatex({\mdseries\slshape data})\index{CtblLegendLatex@\texttt{CtblLegendLatex}}
\label{CtblLegendLatex}
}\hfill{\scriptsize (function)}}\\
\textbf{\indent Returns:\ }
 A String 



 Generates a string representation of the mathematical substitutions \mbox{\texttt{\mdseries\slshape data}}, generated by the undocumented function \texttt{CharacterTableDisplayStringEntryDataDefault} for entries within a character table. 

 }

 

\subsection{\textcolor{Chapter }{GenNameAssocLetterLatex}}
\logpage{[ 2, 1, 4 ]}\nobreak
\hyperdef{L}{X7FF6C91E798C3B3F}{}
{\noindent\textcolor{FuncColor}{$\triangleright$\enspace\texttt{GenNameAssocLetterLatex({\mdseries\slshape s})\index{GenNameAssocLetterLatex@\texttt{GenNameAssocLetterLatex}}
\label{GenNameAssocLetterLatex}
}\hfill{\scriptsize (function)}}\\
\textbf{\indent Returns:\ }
 A String 



 Generates a string representation of the provided letter string \mbox{\texttt{\mdseries\slshape s}} correctly subscripted with a LaTeX math\texttt{\symbol{45}}mode subscript
environment. 

 }

 

\subsection{\textcolor{Chapter }{FactoriseAssocWordLatex}}
\logpage{[ 2, 1, 5 ]}\nobreak
\hyperdef{L}{X7A735F5A7A3BABC9}{}
{\noindent\textcolor{FuncColor}{$\triangleright$\enspace\texttt{FactoriseAssocWordLatex({\mdseries\slshape l, names, tseed})\index{FactoriseAssocWordLatex@\texttt{FactoriseAssocWordLatex}}
\label{FactoriseAssocWordLatex}
}\hfill{\scriptsize (function)}}\\
\textbf{\indent Returns:\ }
 A Factorised String 



 Factorises the string representation of an assoc word in letter representation \mbox{\texttt{\mdseries\slshape l}}, based on the return value from the undocumented function \texttt{FindSubstringPowers}, using the passed list of letters \mbox{\texttt{\mdseries\slshape names}}, and a list of reserved numbers \mbox{\texttt{\mdseries\slshape tseed}} (typically empty for initial calls). 

 This method is essentially a LaTeX\texttt{\symbol{45}}specific implementation
of the the undocumented function \texttt{NiceStringAssocWord}. 

 }

 }

 
\section{\textcolor{Chapter }{Rendering LaTeX Strings}}\label{Chapter_LaTeX_Generation_Section_Rendering_LaTeX_Strings}
\logpage{[ 2, 2, 0 ]}
\hyperdef{L}{X783EC1DF7944F979}{}
{
  

 To aid users using \texttt{Typeset} (\ref{Typeset}), being able to view the results quickly in a variety of
widely\texttt{\symbol{45}}used formats would help streamline the usage of the
package. As such, a number of functions described below have been written to
enable users to render the output LaTeX\texttt{\symbol{45}}renderable snippets
in different fashions. 

 

\subsection{\textcolor{Chapter }{RenderLatex}}
\logpage{[ 2, 2, 1 ]}\nobreak
\hyperdef{L}{X7B9B85188742BD7A}{}
{\noindent\textcolor{FuncColor}{$\triangleright$\enspace\texttt{RenderLatex({\mdseries\slshape str[, options]})\index{RenderLatex@\texttt{RenderLatex}}
\label{RenderLatex}
}\hfill{\scriptsize (function)}}\\


 Renders a given LaTeX string \mbox{\texttt{\mdseries\slshape str}} in a LaTeX environment, providing a visual example of what the string would
look like in a paper. By default, this involved creating a HTML file that
includes the MathJax script, but the \textsf{GAP} option \mbox{\texttt{\mdseries\slshape output}} can be passed to change the rendering method. 

 Currently implemented rendering methods are: 
\begin{itemize}
\item  \texttt{"mathjax"}, calling \texttt{MathJax} (\ref{MathJax}) 
\item  \texttt{"pdflatex"}, calling \texttt{PDFLatex} (\ref{PDFLatex}) 
\item  \texttt{"overleaf"}, calling \texttt{Overleaf} (\ref{Overleaf}) 
\end{itemize}
 

 These functions can be called independently of \texttt{RenderLatex}, which serves only to be a more centralised rendering method for users. 

 }

 

\subsection{\textcolor{Chapter }{PDFLatex}}
\logpage{[ 2, 2, 2 ]}\nobreak
\hyperdef{L}{X87D8B6E98449F0F0}{}
{\noindent\textcolor{FuncColor}{$\triangleright$\enspace\texttt{PDFLatex({\mdseries\slshape str})\index{PDFLatex@\texttt{PDFLatex}}
\label{PDFLatex}
}\hfill{\scriptsize (function)}}\\


 Renders a given LaTeX string \mbox{\texttt{\mdseries\slshape str}} in a new PDF file, specifically via the pdflatex bash tool. 

 }

 

\subsection{\textcolor{Chapter }{MathJax}}
\logpage{[ 2, 2, 3 ]}\nobreak
\hyperdef{L}{X8179071D786C381C}{}
{\noindent\textcolor{FuncColor}{$\triangleright$\enspace\texttt{MathJax({\mdseries\slshape str})\index{MathJax@\texttt{MathJax}}
\label{MathJax}
}\hfill{\scriptsize (function)}}\\


 Renders a given LaTeX string \mbox{\texttt{\mdseries\slshape str}} in a HTML file, making use of the MathJax and TikzJax scripts. 

 }

 

\subsection{\textcolor{Chapter }{Overleaf}}
\logpage{[ 2, 2, 4 ]}\nobreak
\hyperdef{L}{X800F87247AE87FE5}{}
{\noindent\textcolor{FuncColor}{$\triangleright$\enspace\texttt{Overleaf({\mdseries\slshape str})\index{Overleaf@\texttt{Overleaf}}
\label{Overleaf}
}\hfill{\scriptsize (function)}}\\


 Renders a given LaTeX string \mbox{\texttt{\mdseries\slshape str}} in a new Overleaf project, specifically via a URL\texttt{\symbol{45}}encoded
snippet. 

 }

 

\subsection{\textcolor{Chapter }{URIEncodeComponent}}
\logpage{[ 2, 2, 5 ]}\nobreak
\hyperdef{L}{X824C38568286ECBB}{}
{\noindent\textcolor{FuncColor}{$\triangleright$\enspace\texttt{URIEncodeComponent({\mdseries\slshape raw})\index{URIEncodeComponent@\texttt{URIEncodeComponent}}
\label{URIEncodeComponent}
}\hfill{\scriptsize (function)}}\\
\textbf{\indent Returns:\ }
 A Percent\texttt{\symbol{45}}Encoded String 



 Replaces reserved characters within a URI component \mbox{\texttt{\mdseries\slshape raw}} as per RFC\texttt{\symbol{45}}3986. 

 }

 

\subsection{\textcolor{Chapter }{NeedsLatexMathMode}}
\logpage{[ 2, 2, 6 ]}\nobreak
\hyperdef{L}{X8628E043862C318A}{}
{\noindent\textcolor{FuncColor}{$\triangleright$\enspace\texttt{NeedsLatexMathMode({\mdseries\slshape raw})\index{NeedsLatexMathMode@\texttt{NeedsLatexMathMode}}
\label{NeedsLatexMathMode}
}\hfill{\scriptsize (function)}}\\
\textbf{\indent Returns:\ }
 A Boolean 



 Determines if the LaTeX snippet \mbox{\texttt{\mdseries\slshape raw}} requires LaTeX's math mode to be rendered correctly. 

 }

 

\subsection{\textcolor{Chapter }{DEFAULT{\textunderscore}LATEX{\textunderscore}PREAMBLE}}
\logpage{[ 2, 2, 7 ]}\nobreak
\hyperdef{L}{X853D2C5F85183D3C}{}
{\noindent\textcolor{FuncColor}{$\triangleright$\enspace\texttt{DEFAULT{\textunderscore}LATEX{\textunderscore}PREAMBLE\index{DEFAULT{\textunderscore}LATEX{\textunderscore}PREAMBLE@\texttt{DEF}\-\texttt{A}\-\texttt{U}\-\texttt{L}\-\texttt{T{\textunderscore}}\-\texttt{L}\-\texttt{A}\-\texttt{T}\-\texttt{E}\-\texttt{X{\textunderscore}}\-\texttt{P}\-\texttt{R}\-\texttt{E}\-\texttt{A}\-\texttt{MBLE}}
\label{DEFAULTuScoreLATEXuScorePREAMBLE}
}\hfill{\scriptsize (global variable)}}\\


 Default LaTeX preamble string used for creating compilable \texttt{.tex} files from LaTeX snippets. }

 

\subsection{\textcolor{Chapter }{DEFAULT{\textunderscore}MATHJAX{\textunderscore}TAGS}}
\logpage{[ 2, 2, 8 ]}\nobreak
\hyperdef{L}{X83CE796886A570A5}{}
{\noindent\textcolor{FuncColor}{$\triangleright$\enspace\texttt{DEFAULT{\textunderscore}MATHJAX{\textunderscore}TAGS\index{DEFAULT{\textunderscore}MATHJAX{\textunderscore}TAGS@\texttt{DEF}\-\texttt{A}\-\texttt{U}\-\texttt{L}\-\texttt{T{\textunderscore}}\-\texttt{M}\-\texttt{A}\-\texttt{T}\-\texttt{H}\-\texttt{J}\-\texttt{A}\-\texttt{X{\textunderscore}}\-\texttt{TAGS}}
\label{DEFAULTuScoreMATHJAXuScoreTAGS}
}\hfill{\scriptsize (global variable)}}\\


 Default HTML document and head tags used to create HTML files using MathJax to
render LaTeX snippets. }

 

\subsection{\textcolor{Chapter }{ALWAYS{\textunderscore}UNESCAPED{\textunderscore}CHARS}}
\logpage{[ 2, 2, 9 ]}\nobreak
\hyperdef{L}{X789A6A708289DE47}{}
{\noindent\textcolor{FuncColor}{$\triangleright$\enspace\texttt{ALWAYS{\textunderscore}UNESCAPED{\textunderscore}CHARS\index{ALWAYS{\textunderscore}UNESCAPED{\textunderscore}CHARS@\texttt{ALW}\-\texttt{A}\-\texttt{Y}\-\texttt{S{\textunderscore}}\-\texttt{U}\-\texttt{N}\-\texttt{E}\-\texttt{S}\-\texttt{C}\-\texttt{A}\-\texttt{P}\-\texttt{E}\-\texttt{D{\textunderscore}}\-\texttt{C}\-\texttt{HARS}}
\label{ALWAYSuScoreUNESCAPEDuScoreCHARS}
}\hfill{\scriptsize (global variable)}}\\


 String containing all of the characters that do not need to be
percent\texttt{\symbol{45}}encoded within URI components, as per
RFC\texttt{\symbol{45}}3986. }

 }

 
\section{\textcolor{Chapter }{Digraphs Integration}}\label{Chapter_LaTeX_Generation_Section_Digraphs_Integration}
\logpage{[ 2, 3, 0 ]}
\hyperdef{L}{X7EF74AB68325E735}{}
{
  

 \textsf{digraphs} is a powerful, widely\texttt{\symbol{45}}used packages, implementing helpful
functionality to work with directed graphs amongst other objects. Due to it's
popularity, and as a way to demonstrate how \textsf{typeset} can be integrated with external packages, the following functions have been
implemented to allow directed graphs to be converted into LaTeX
representations. 

 It should be noted that the conversion implemented here does use the output
from \texttt{DotDigraph} (\textbf{Digraphs: DotDigraph}), which generates the DOT string representing a digraph. This is then used to
either convert it to a tikz representation via the
command\texttt{\symbol{45}}line tool dot2tex (using the \textsf{GAP} option \texttt{digraphOut := "dot2tex"}), or simply wrapping it up in a dot2tex environment provided by the LaTeX
package dot2texi which will compile the wrapped DOT input into tikz during
compilation of the LaTeX file itself. 

 While another method could be written to convert the internal representation
of a directed graph directly into a tikzpicture, this would likely be
incredibly convoluted and difficult, and may present numerous problems with
edge positioning. Therefore, relying on dot2tex was chosen as the best
approach. 

 

\subsection{\textcolor{Chapter }{DEFAULT{\textunderscore}DOT2TEX{\textunderscore}OPTIONS}}
\logpage{[ 2, 3, 1 ]}\nobreak
\hyperdef{L}{X7AD116C5832E82C1}{}
{\noindent\textcolor{FuncColor}{$\triangleright$\enspace\texttt{DEFAULT{\textunderscore}DOT2TEX{\textunderscore}OPTIONS\index{DEFAULT{\textunderscore}DOT2TEX{\textunderscore}OPTIONS@\texttt{DEF}\-\texttt{A}\-\texttt{U}\-\texttt{L}\-\texttt{T{\textunderscore}}\-\texttt{D}\-\texttt{O}\-\texttt{T2}\-\texttt{T}\-\texttt{E}\-\texttt{X{\textunderscore}}\-\texttt{O}\-\texttt{P}\-\texttt{T}\-\texttt{IONS}}
\label{DEFAULTuScoreDOT2TEXuScoreOPTIONS}
}\hfill{\scriptsize (global variable)}}\\


 Default command\texttt{\symbol{45}}line options passed to the \texttt{dot2tex} executable to convert dot strings. }

 

\subsection{\textcolor{Chapter }{Dot2Tex}}
\logpage{[ 2, 3, 2 ]}\nobreak
\hyperdef{L}{X86E66F1C7816D90D}{}
{\noindent\textcolor{FuncColor}{$\triangleright$\enspace\texttt{Dot2Tex({\mdseries\slshape obj})\index{Dot2Tex@\texttt{Dot2Tex}}
\label{Dot2Tex}
}\hfill{\scriptsize (function)}}\\
\textbf{\indent Returns:\ }
 A Tikz String 



 Executes \texttt{dot2tex} on the dot string representing a given digraph object \mbox{\texttt{\mdseries\slshape obj}}. 

 }

 }

 

 }

   
\chapter{\textcolor{Chapter }{Structure Descriptions}}\label{Chapter_Structure_Descriptions}
\logpage{[ 3, 0, 0 ]}
\hyperdef{L}{X87BF1B887C91CA2E}{}
{
  

 
\section{\textcolor{Chapter }{Typesetting Structure Descriptions of Groups}}\label{Chapter_Structure_Descriptions_Section_Typesetting_Structure_Descriptions_of_Groups}
\logpage{[ 3, 1, 0 ]}
\hyperdef{L}{X7DC1724784ACBEF0}{}
{
  

 

\subsection{\textcolor{Chapter }{TypesetStructureDescription}}
\logpage{[ 3, 1, 1 ]}\nobreak
\hyperdef{L}{X807F54B97A98F647}{}
{\noindent\textcolor{FuncColor}{$\triangleright$\enspace\texttt{TypesetStructureDescription({\mdseries\slshape desc})\index{TypesetStructureDescription@\texttt{TypesetStructureDescription}}
\label{TypesetStructureDescription}
}\hfill{\scriptsize (function)}}\\
\textbf{\indent Returns:\ }
 A String 



 Generates a typesettable representation equivalent to a given structure
description \mbox{\texttt{\mdseries\slshape desc}} of a group. Said structure descriptions can be calculated via \texttt{StructureDescription} (\textbf{Reference: StructureDescription}). 

 }

 

\subsection{\textcolor{Chapter }{LatexStructureDescription}}
\logpage{[ 3, 1, 2 ]}\nobreak
\hyperdef{L}{X7A6316BB879E51FD}{}
{\noindent\textcolor{FuncColor}{$\triangleright$\enspace\texttt{LatexStructureDescription({\mdseries\slshape desc})\index{LatexStructureDescription@\texttt{LatexStructureDescription}}
\label{LatexStructureDescription}
}\hfill{\scriptsize (function)}}\\
\textbf{\indent Returns:\ }
 A String 



 Generates a LaTeX representation equivalent to a given structure description \mbox{\texttt{\mdseries\slshape desc}}. Called by \texttt{TypesetStructureDescription} (\ref{TypesetStructureDescription}) as the default markup language for structure description typesetting. 

 }

 

\subsection{\textcolor{Chapter }{ConcatStructDescOperands}}
\logpage{[ 3, 1, 3 ]}\nobreak
\hyperdef{L}{X86362669792C2F28}{}
{\noindent\textcolor{FuncColor}{$\triangleright$\enspace\texttt{ConcatStructDescOperands({\mdseries\slshape desc, sep, newsep})\index{ConcatStructDescOperands@\texttt{ConcatStructDescOperands}}
\label{ConcatStructDescOperands}
}\hfill{\scriptsize (function)}}\\
\textbf{\indent Returns:\ }
 A String 



 Concatenates the tokens parsed from splitting the string \mbox{\texttt{\mdseries\slshape desc}} with the provided separator string \mbox{\texttt{\mdseries\slshape sep}}. It will then process the tokens with \texttt{LatexStructureDescription} (\ref{LatexStructureDescription}), and concatenate the results with the given new separator \mbox{\texttt{\mdseries\slshape newsep}}. 

 }

 }

 

 While \texttt{StructureDescription} (\textbf{Reference: StructureDescription}) provides powerful functionality in calculating the structure description of an
input group, the returned string is simply raw text, and does not look good
when rendered in typical typesetting environment. 

 To improve this, the following functions have been written to help convert the
structure description strings into typesettable representations, that use
language\texttt{\symbol{45}}specific features to ensure that displaying
structure descriptions aesthetically is easy and efficient. 

 }

   
\chapter{\textcolor{Chapter }{MathML Example Generation}}\label{Chapter_MathML_Example_Generation}
\logpage{[ 4, 0, 0 ]}
\hyperdef{L}{X80E5E58082AB0716}{}
{
  
\section{\textcolor{Chapter }{MathML Generation Functions for \textsf{GAP} Objects}}\label{Chapter_MathML_Example_Generation_Section_MathML_Generation_Functions_for_GAP_Objects}
\logpage{[ 4, 1, 0 ]}
\hyperdef{L}{X781F62C37D928E4C}{}
{
  

 This section describes a bare\texttt{\symbol{45}}bones implementation of the
framework for generating MathML representations. It is intended to be used as
a reference for implementing the framework, and as a starting point for
implementing the framework for other languages \texttt{\symbol{45}} not as a
fully functional implementation of the framework. 

\subsection{\textcolor{Chapter }{GenMathmlTmpl (for IsObject)}}
\logpage{[ 4, 1, 1 ]}\nobreak
\hyperdef{L}{X7E9CF77579FA1393}{}
{\noindent\textcolor{FuncColor}{$\triangleright$\enspace\texttt{GenMathmlTmpl({\mdseries\slshape obj})\index{GenMathmlTmpl@\texttt{GenMathmlTmpl}!for IsObject}
\label{GenMathmlTmpl:for IsObject}
}\hfill{\scriptsize (operation)}}\\
\textbf{\indent Returns:\ }
 An Unpopulated MathML Format String 



 Generates a format string that represents the structural definition of the
given \textsf{GAP} object \mbox{\texttt{\mdseries\slshape obj}} in MathML. It contains no parameter values, and will need to be populated with
the arguments representing the semantic values of the object, generated via \texttt{GenArgs} (\ref{GenArgs:for IsObject}), before it can be rendered in a MathML environment. 

 }

 }

 }

 \def\indexname{Index\logpage{[ "Ind", 0, 0 ]}
\hyperdef{L}{X83A0356F839C696F}{}
}

\cleardoublepage
\phantomsection
\addcontentsline{toc}{chapter}{Index}


\printindex

\immediate\write\pagenrlog{["Ind", 0, 0], \arabic{page},}
\newpage
\immediate\write\pagenrlog{["End"], \arabic{page}];}
\immediate\closeout\pagenrlog
\end{document}
